% Part: history
% Chapter: biographies 
% Section: kurt-goedel

\documentclass[../../../include/open-logic-section]{subfiles}

\begin{document}

\olfileid[pt]{his}{bio}{god}

\olsection{Kurt G\"odel}

\olphoto{goedel-kurt}{Kurt G{\"o}del}

Kurt G{\"o}del (\textsc{ger}-dle) nasceu a 28 de abril de 1906 em
Br{\"u}nn, no império austro-húngaro (hoje Brno, na República Checa).
Devido ao seu espírito inquisitivo e brilhante, o jovem Kurtele era
muitas vezes chamado ``Der kleine Herr Warum'' (o Pequeno Senhor
Porquê) pela família. Destacou-se nos estudos desde a escola
primária, tendo obtido menos do que a nota mais alta apenas em
matemática. G{\"o}del faltava frequentemente à escola por causa da
frágil saúde e estava isento de educação física. Foi diagnosticado
com febre reumática durante a infância. Ao longo da vida, acreditou
que isto tinha afetado permanentemente o seu coração, apesar dos
pareceres médicos em contrário.

G{\"o}del começou a estudar na Universidade de Viena em 1924 e
concluiu o doutoramento em~1929. Inicialmente pretendia estudar
física, mas os seus interesses deslocaram-se rapidamente para a
matemática e sobretudo para a lógica, em parte devido à influência do
filósofo Rudolf Carnap. A sua dissertação, orientada por Hans Hahn,
provava o teorema de completude da lógica de predicados de primeira
ordem com identidade \citep{Godel1929}. Apenas um ano depois, obteve
os seus resultados mais famosos---o primeiro e o segundo teorema da
incompletude (publicados em \citealt{Godel1931}). Durante o período
em Viena, G{\"o}del participou intensamente no Círculo de Viena, um
grupo de filósofos de orientação científica que incluía Carnap, cujo
trabalho foi especialmente influenciado pelos resultados de G{\"o}del.

Em 1938, G\"odel casou com Adele Nimbursky. Os pais não ficaram
contentes: não só ela era seis anos mais velha e já divorciada, como
trabalhava como bailarina num clube noturno.
As pressões sociais não afetaram G{\"o}del, contudo,
e permaneceram felizmente casados até à sua morte.

Depois de a Alemanha Nazi ter anexado a Áustria em 1938, G{\"o}del e
Adele emigraram para os Estados Unidos, onde ele assumiu um lugar no
Institute for Advanced Study em Princeton, Nova Jérsia. Apesar da
introversão e do caráter excêntrico, o período de G{\"o}del em
Princeton foi colaborativo e frutuoso.  Publicou ensaios em teoria dos
conjuntos, filosofia e física. Notavelmente, estabeleceu uma amizade
particularmente forte com o colega no IAS, Albert Einstein.

Nos últimos anos, a saúde mental de G{\"o}del deteriorou-se. A
internação da mulher em 1977 fez com que ela já não pudesse cozinhar-lhe
as refeições. Tendo sofrido de problemas de saúde mental ao longo da
vida, sucumbiu à paranoia. Com medo mortal de ser envenenado,
G{\"o}del recusou-se a comer. Morreu de inanição a 14 de janeiro de
1978, em Princeton.

\begin{reading}
Para uma biografia completa da vida de G{\"o}del, ver
\citet{Dawson1997}. Para outros textos biográficos, bem como ensaios
sobre as contribuições de G{\"o}del para a lógica e a filosofia, ver
\citet{Wang1990}, \citet{Baaz2011}, \citet{Takeuti2003}, e
\citet{Sigmund2007}.

A tese de doutoramento de G{\"o}del está disponível no alemão
original \citep{Godel1929}.  O texto original dos teoremas da
incompletude é \citep{Godel1931}. Todas as publicações e textos não
publicados de G\"odel, bem como uma seleção de correspondência, estão
disponíveis em inglês nos seus \emph{Collected Papers}
\citet{Godel1986,Godel1990}.

Para um tratamento detalhado dos teoremas da incompletude de
G{\"o}del, ver \citet{Smith2013}. Para uma discussão informal e
filosófica dos teoremas de G{\"o}del, ver o podcast de Mark
Linsenmayer \citep{Linsenmayer2014}.
\end{reading}

\end{document}
